% ============================================================================
% CS455 Term Project — Knowledge Distillation of Mathematical Reasoning
% Ali Kumral, Revna Demirkale
%
% OVERLEAF INSTRUCTIONS:
%   1. New Project -> Upload Project (or paste this file).
%   2. Upload the two figures from results/figures/ on your Drive:
%        F1_difficulty_accuracy.png   F2_generation_lengths.png
%      into the same folder (or a figures/ subfolder; adjust \graphicspath).
%   3. Compile with pdfLaTeX.
% ============================================================================
\documentclass[11pt]{article}

\usepackage[utf8]{inputenc}
\usepackage[T1]{fontenc}
\usepackage[margin=1in]{geometry}
\usepackage{booktabs}
\usepackage{graphicx}
\usepackage{amsmath,amssymb}
\usepackage{microtype}
\usepackage[hidelinks]{hyperref}
\usepackage{caption}
\usepackage{authblk}

\graphicspath{{figures/}{./}}

\title{\textbf{Knowledge Distillation of Mathematical Reasoning:\\
Transferring Chain-of-Thought from Llama-3.1-8B to Llama-3.2-1B}}
\author[]{Ali Kumral (30586) \quad Revna Demirkale (32056)}
\affil[]{CS 455 --- Large Language Models, Spring 2025/2026 \quad Category (D): Efficiency / Optimization}
\date{}

\begin{document}
\maketitle

\begin{abstract}
\noindent
Large language models for multi-step reasoning are expensive to serve, motivating
the transfer of reasoning ability from large \emph{teacher} models into small
\emph{student} models. We study sequence-level chain-of-thought (CoT) distillation
from \textbf{Llama-3.1-8B-Instruct} into \textbf{Llama-3.2-1B-Instruct} on
grade-school math, using a tight, fully reproducible recipe that runs on Google
Colab. We generate teacher CoT traces over the GSM8K training set, filter them by
rejection sampling against gold answers (81.5\% acceptance), and fine-tune the
student with QLoRA under five controlled conditions. Distillation works and is
statistically overwhelming: the distilled 1B student improves from \textbf{38.1\%}
to \textbf{46.9\%} on GSM8K test ($+8.9$ pp, McNemar $p<10^{-9}$), and crucially
\textbf{outperforms a student trained on GSM8K's own gold solutions} by $+13$ pp
($p<10^{-17}$), showing that teacher-generated reasoning carries signal beyond the
labeled data. We also report two robust \emph{negative} findings: (1) fine-tuning
on gold CoT \emph{degrades} accuracy below the zero-shot baseline, and (2) using
three teacher traces per problem is \emph{significantly worse} than one (a
$\sim$6$\sigma$ effect across three seeds), contrary to our hypothesis. Gains do
not transfer out of distribution: accuracy on a held-out MATH subset drops relative
to zero-shot. Finally, the practical efficiency advantage of the 1B student is
\textbf{memory} ($5.2\times$ lower peak VRAM), not latency, because the distilled
student inherits the teacher's verbosity. We release code, the cleaned trace
dataset, and the trained adapter.
\end{abstract}

\section{Introduction}
Deploying LLMs for reasoning tasks---math, code, structured planning---is
bottlenecked by the cost and latency of large models. A 1B model is roughly
$8\times$ cheaper to serve than an 8B model and dramatically cheaper than a 70B
model; if a small student can be brought close to a larger teacher's quality on a
target task, edge inference, on-device assistants, and high-throughput batch jobs
all become viable. Chain-of-thought (CoT) distillation---training a small model on
the step-by-step reasoning of a larger one---has been shown to yield large gains,
but published recipes are often computationally heavy or rely on closed teacher
APIs.

We ask a focused, practical question: \emph{how far can a 1B student be pushed on
grade-school math reasoning by distilling from a small, open 8B teacher, under a
single-digit-GPU-hour budget, with rigorous baselines?} Rather than chasing a
single headline number, we build five tightly controlled conditions so that every
comparison isolates one variable, and we report honest negative results where they
occur.

\paragraph{Contributions.}
\begin{enumerate}
  \item A tight, fully open, reproducible CoT-distillation recipe (8B teacher, not
  70B; QLoRA; runs end-to-end on Colab) with resumable pipeline stages and
  released artifacts.
  \item Rigorous evidence that distillation works and that the teacher signal
  exceeds gold labels (H1, H2), confirmed with bootstrap confidence intervals and
  exact McNemar tests across three random seeds.
  \item Two robust negative findings: training on gold CoT \emph{hurts} relative to
  zero-shot, and \emph{more} teacher traces per problem (3 vs.\ 1) is significantly
  \emph{worse}---a $\sim$6$\sigma$ effect.
  \item An efficiency analysis showing the real lever is memory, not latency, due
  to verbosity inheritance, plus per-difficulty and qualitative error analyses and
  an out-of-distribution (MATH) evaluation showing no transfer.
\end{enumerate}

\section{Related Work}
\paragraph{Knowledge distillation.} Hinton et al.~\cite{hinton2015} introduced
distillation via soft-label matching. For autoregressive LLMs, storing teacher
logits is impractical, so \emph{sequence-level} distillation---training the student
on text generated by the teacher---has become standard. MiniLLM~\cite{minillm} and
DistiLLM~\cite{distillm} refine the objective and data construction and report
sequence-level distillation as competitive. The DeepSeek-R1 report~\cite{deepseekr1}
shows that fine-tuning small models on long CoT traces from a strong teacher
transfers reasoning effectively.

\paragraph{CoT distillation for reasoning.} Magister et al.~\cite{magister2023}
show that fine-tuning small models on teacher CoT improves multi-step arithmetic,
and that \emph{rejection sampling} against gold answers improves trace quality---the
recipe we follow.

\paragraph{Efficient fine-tuning.} LoRA~\cite{lora} freezes the base model and
learns low-rank adapters; QLoRA~\cite{qlora} adds 4-bit NF4 quantization of the
frozen base, enabling fine-tuning of billion-parameter models on a single GPU. We
use QLoRA throughout.

\paragraph{Benchmarks.} GSM8K~\cite{gsm8k} is the canonical grade-school math
word-problem benchmark with gold CoT; MATH~\cite{math} contains harder
competition-level problems and serves as our out-of-distribution test.

Our work differs from prior distillation papers in deliberately staying small and
open (8B teacher, 1B student, single-GPU budget) and in emphasizing controlled
baselines and honest negative results over a single benchmark number.

\section{Method}
\subsection{Overview}
Our pipeline has four stages: (1) teacher CoT generation, (2) rejection sampling
and dataset construction, (3) QLoRA fine-tuning of the student, and (4) evaluation.
All logic lives in a small installable package; each stage reads and writes
versioned JSONL artifacts so that long runs on Colab are resumable across
disconnects.

\subsection{Comparison conditions}
We train and evaluate five conditions (Table~\ref{tab:conditions}), designed so
that each pairwise comparison isolates a single factor. Conditions ii--iv are
trained with three random seeds; conditions i and v are deterministic (greedy
decoding). Gold and teacher training targets are formatted identically (same system
prompt, same final-answer line), so the \emph{only} difference between conditions
ii and iii is the \emph{content} of the reasoning, making the H2 comparison fair.

\begin{table}[t]
\centering
\caption{The five comparison conditions.}
\label{tab:conditions}
\begin{tabular}{clccl}
\toprule
\# & Condition & Trained? & Training data & Isolates \\
\midrule
i   & Zero-shot 1B      & no  & ---                       & lower bound \\
ii  & Gold CoT          & yes & GSM8K gold solutions      & labels alone \\
iii & Teacher-1         & yes & 1 teacher trace / problem & distillation (H1) \\
iv  & Teacher-3         & yes & $\le$3 traces / problem   & more traces (H3) \\
v   & 8B teacher        & no  & ---                       & upper bound \\
\bottomrule
\end{tabular}
\end{table}

\subsection{Teacher CoT generation}
We run Llama-3.1-8B-Instruct in 4-bit NF4 (double quantization, bf16 compute) over
all 7{,}473 GSM8K training problems. To obtain diverse traces while controlling
VRAM, we sample in \textbf{three independent passes} (one trace per problem per
pass, temperature 0.7, top-$p$ 0.9, max 256 new tokens) rather than drawing $k$
samples in a single decoding call. Each pass uses a different seed and appends to a
shared pool; generation is resumable at the (problem, sample) level. This produced
22{,}449 raw traces.

\subsection{Rejection sampling and dataset construction}
We parse the final numerical answer from each trace and keep only traces matching
the gold label. \textbf{18{,}287 / 22{,}449 traces were accepted (81.5\%)},
covering 6{,}922 of 7{,}473 problems with $\ge$1 correct trace and 5{,}081 with
$\ge$3. We build three SFT datasets: \emph{Gold} (7{,}473 examples), \emph{Teacher-1}
(6{,}922; one trace/problem), and \emph{Teacher-3} (18{,}263; up to three/problem).

\subsection{QLoRA fine-tuning}
The student (Llama-3.2-1B-Instruct) is loaded in 4-bit NF4; we attach LoRA adapters
(rank 16, $\alpha=32$, dropout 0.05) on all attention and MLP projections. We train
for 3 epochs at learning rate $2\times10^{-4}$ (cosine schedule, 3\% warmup),
effective batch size 16, bf16, gradient checkpointing, and paged 8-bit AdamW, with
maximum sequence length 1{,}024. Training one condition takes 18--49 minutes on an
A100.

\paragraph{Implementation note (loss masking).} We intended \emph{completion-only}
loss masking (loss on the assistant reasoning, not the prompt). The TRL version
available on Colab had removed the collator we targeted and gated its built-in
masking behind chat-template markers absent from Llama-3, so all conditions were
trained with \textbf{full-sequence loss} instead. Because this choice is applied
identically to every condition, it does not affect the \emph{relative} comparisons
our hypotheses test; completion-only masking is left as future work
(Section~\ref{sec:limitations}).

\subsection{Evaluation protocol}
We evaluate with greedy decoding (deterministic) and extract the final numerical
answer using a fixed marker plus fallbacks ($\backslash$boxed\{\}, \texttt{\#\#\#\#},
last number) with numeric normalization. The primary metric is GSM8K test
exact-match accuracy (1{,}319 problems). Secondary metrics: MATH out-of-distribution
accuracy (a frozen, level-stratified subset of 167 problems, ``MATH-167''),
single-example latency (batch size 1), peak inference VRAM, and generation length.
We report 95\% bootstrap confidence intervals over test items and, for paired
comparisons, exact McNemar tests; trained conditions additionally report mean
$\pm$ std over three seeds.

\section{Experimental Setup}
\begin{description}
  \item[Models.] Teacher \texttt{meta-llama/Llama-3.1-8B-Instruct}; student
  \texttt{meta-llama/Llama-3.2-1B-Instruct}. Both loaded in 4-bit NF4 with double
  quantization and bf16 compute.
  \item[Data.] \texttt{openai/gsm8k} (config \texttt{main}): 7{,}473 train /
  1{,}319 test. Out-of-distribution: a frozen, level-stratified subset of
  \texttt{qwedsacf/competition\_math} (167 problems, ``MATH-167'').
  \item[Hardware.] Google Colab, NVIDIA A100 for generation, training, and most
  evaluation. (The proposal specified latency on a T4; we report A100 latency and
  treat absolute values as hardware-specific---the cross-model \emph{ratio} is the
  meaningful quantity.)
  \item[Reproducibility.] All hyperparameters live in version-controlled YAML
  configs; every run records its seed and configuration; pipeline stages are
  resumable.
\end{description}

\section{Results}
\subsection{Main results}
Table~\ref{tab:main} reports accuracy with 95\% bootstrap CIs (seed 0) and seed
mean $\pm$ std for trained conditions; Table~\ref{tab:mcnemar} reports paired
significance.

\begin{table}[t]
\centering
\caption{Main results. Accuracy (\%); CIs are 95\% bootstrap (seed 0).
Teacher-1/Teacher-3 show mean\,$\pm$\,std over 3 seeds.}
\label{tab:main}
\begin{tabular}{lcccc}
\toprule
Condition & GSM8K & GSM8K 95\% CI & MATH & MATH 95\% CI \\
\midrule
Zero-shot 1B        & 38.06            & [35.5, 40.6] & 33.53            & [26.4, 40.7] \\
Gold CoT            & 33.74            & [31.2, 36.3] & 20.96            & [15.0, 27.5] \\
\textbf{Teacher-1}  & \textbf{46.93\,$\pm$\,0.11} & [44.1, 49.5] & 29.14\,$\pm$\,1.98 & [20.4, 33.5] \\
Teacher-3           & 44.33\,$\pm$\,0.77 & [40.7, 46.1] & 26.15\,$\pm$\,0.56 & [20.4, 33.5] \\
8B teacher          & 82.87            & [80.8, 84.9] & 53.89            & [46.1, 61.7] \\
\bottomrule
\end{tabular}
\end{table}

\begin{table}[t]
\centering
\caption{Exact McNemar tests, paired on the 1{,}319 GSM8K test items.
``A-only / B-only'' = problems one model got right that the other missed.}
\label{tab:mcnemar}
\begin{tabular}{llcl}
\toprule
Comparison & A-only / B-only & $p$ & Verdict \\
\midrule
H1: Teacher-1 vs Zero-shot & 236 / 121 & $1.2\times10^{-9}$  & distillation works *** \\
H2: Teacher-1 vs Gold      & 290 / 118 & $8.2\times10^{-18}$ & teacher $\gg$ gold *** \\
H3: Teacher-1 vs Teacher-3 & 199 / 154 & $0.019$             & one trace \emph{beats} three * \\
Gold vs Zero-shot          & 172 / 229 & $0.005$             & gold \emph{hurts} ** \\
\bottomrule
\end{tabular}
\end{table}

\paragraph{Findings.}
\begin{itemize}
  \item \textbf{H1 (distillation works) --- confirmed, overwhelmingly.} Teacher-1
  improves over zero-shot by $+8.9$ pp ($p<10^{-9}$); a seed std of $0.0011$ shows
  the result is extremely stable.
  \item \textbf{H2 (teacher $>$ gold) --- confirmed, overwhelmingly.} Teacher-1
  beats Gold by $+13.2$ pp ($p<10^{-17}$): teacher-generated reasoning provides
  signal beyond the labeled data.
  \item \textbf{H3 (more traces help) --- reversed and robust.} Teacher-1
  ($46.93\pm0.11$) significantly beats Teacher-3 ($44.33\pm0.77$); the $+2.6$ pp gap
  is $\approx5.9\times$ the pooled seed std ($p=0.019$).
  \item \textbf{Gold fine-tuning hurts.} Condition ii falls $4.3$ pp \emph{below}
  zero-shot ($p=0.005$).
\end{itemize}

\subsection{Accuracy by difficulty}
We bucket the 1{,}319 test problems into easy/medium/hard tertiles by gold-CoT
length ($\le$36 / 37--57 / $>$57 words; 452/431/436 problems). Distillation helps
at all levels, with large absolute gains on easy ($+11.7$ pp) and hard ($+9.4$ pp)
problems (Table~\ref{tab:difficulty}, Figure~\ref{fig:f1}). The largest remaining
student--teacher gap is on hard problems (32.6 vs.\ 72.7), confirming the 1B
student still cannot match reasoning that tracks several quantities at once. Gold
training notably degrades medium and hard problems ($-7.0$ and $-5.5$ pp).

\begin{table}[t]
\centering
\caption{GSM8K accuracy (\%) by problem difficulty (seed 0).}
\label{tab:difficulty}
\begin{tabular}{lccc}
\toprule
Condition & Easy ($\le$36) & Medium (37--57) & Hard ($>$57) \\
\midrule
Zero-shot 1B & 50.9 & 39.7 & 23.2 \\
Gold         & 50.2 & 32.7 & 17.7 \\
Teacher-1    & 62.6 & 44.6 & 32.6 \\
Teacher-3    & 58.4 & 41.8 & 29.4 \\
8B teacher   & 88.9 & 86.8 & 72.7 \\
\bottomrule
\end{tabular}
\end{table}

\begin{figure}[t]
\centering
\includegraphics[width=0.85\linewidth]{F1_difficulty_accuracy.png}
\caption{GSM8K accuracy by difficulty tertile. Distillation (Teacher-1/3) lifts all
buckets; Gold degrades medium/hard; the teacher gap is largest on hard problems.}
\label{fig:f1}
\end{figure}

\subsection{Generation length and the verbosity tax}
Median generated tokens: Gold 68, Zero-shot 156, Teacher-1 172, Teacher-3 174, 8B
teacher 181 (Figure~\ref{fig:f2}). Two observations: (1) gold training
\emph{collapses} reasoning length to 68 tokens---the model learned to answer
tersely rather than reason, explaining why it underperforms even zero-shot; (2) the
distilled students \emph{inherit the teacher's verbosity} ($\approx$172 vs.\ 181
tokens), nearly matching the 8B model's output length.

\begin{figure}[t]
\centering
\includegraphics[width=\linewidth]{F2_generation_lengths.png}
\caption{Generation-length distributions. Gold collapses to short answers; the
distilled students inherit the teacher's verbosity.}
\label{fig:f2}
\end{figure}

\subsection{Efficiency}
At batch size 1, the distilled 1B student (12.4 s) is barely faster than the 8B
teacher (13.7 s), because verbosity inheritance roughly doubles its output length
relative to zero-shot (6.2 s). The genuine advantage is \textbf{memory: the student
uses $5.2\times$ less peak VRAM} (1.13 vs.\ 5.87 GB; Table~\ref{tab:efficiency}),
which translates into far higher achievable concurrency on fixed hardware.
Distillation buys \emph{density}, not single-query speed.

\begin{table}[t]
\centering
\caption{Inference efficiency (batch size 1).}
\label{tab:efficiency}
\begin{tabular}{lcc}
\toprule
Condition & Latency (s/problem) & Peak VRAM (GB) \\
\midrule
Zero-shot 1B        & 6.17  & 1.09 \\
Teacher-1 (distilled) & 12.39 & 1.13 \\
Teacher-3           & 12.91 & 1.13 \\
8B teacher          & 13.73 & 5.87 \\
\bottomrule
\end{tabular}
\end{table}

\subsection{Out-of-distribution (MATH)}
Distillation does \emph{not} transfer to MATH. Teacher-1 scores $29.14\pm1.98$,
\emph{below} zero-shot's 33.53 and far below the teacher's 53.89. Teacher-1 and
Teacher-3 are statistically indistinguishable on MATH, suggesting MATH performance
is driven by base-model pretraining rather than our GSM8K fine-tuning. A manual
audit of 15 MATH cases marked incorrect found only 1 false negative (LaTeX fraction
formatting), so we treat MATH accuracies as slight lower bounds; the absence of OOD
transfer is robust.

\section{Discussion}
\paragraph{Why does gold fine-tuning hurt?} GSM8K's gold solutions are terse and
stylistically different from the student's natural output. Training on them with
full-sequence loss pushed the student toward short, conclusory answers (median 68
tokens), suppressing the step-by-step reasoning it needs. The teacher's traces, in
the same Llama chat style the student already favors, reinforce rather than fight
the student's prior.

\paragraph{Why does one teacher trace beat three?} Two non-exclusive explanations.
First, with three diverse traces per problem and full-sequence loss, the student
fits three different surface reasoning styles for the same answer, which may blur
its decoding distribution. Second, Teacher-3 reaches a much lower training loss
($\approx$0.22 vs.\ $\approx$0.38) without improving test accuracy---a classic
over-fitting signature. The verbosity statistics rule out a length explanation
(Teacher-1 and Teacher-3 are nearly identical in output length), pointing to
training dynamics rather than data style.

\paragraph{Qualitative trace study.} Of the 1{,}319 test items, Teacher-1 corrected
236 that zero-shot missed while losing 121, with 381 both-correct and 581
both-wrong. Side-by-side cases illustrate the mechanism. On a multi-step
``work-backwards'' problem, zero-shot applied operations in the wrong order while
Teacher-1 introduced a variable and solved the equation systematically---evidence
that the student learned a \emph{reasoning strategy}, not just answers. A
representative failure shows the limit: on an overtime-pay problem Teacher-1 set the
problem up correctly but multiplied the overtime rate by all 45 hours instead of
the 5 overtime hours---a single arithmetic slip inside otherwise correct reasoning.
On easy problems both models already succeed and differ only in phrasing.
(Transcripts in Appendix~\ref{app:qual}.)

\paragraph{Gap closed.} The distilled student closes roughly 20\% of the
zero-shot-to-teacher accuracy gap (8.9 of 44.8 pp) at one-fifth the memory
footprint---a meaningful, honest result at this scale.

\section{Limitations}
\label{sec:limitations}
\begin{itemize}
  \item \textbf{Full-sequence loss.} A TRL API change forced full-sequence rather
  than completion-only loss. It is consistent across conditions (comparisons remain
  fair) but may have amplified the gold-CoT degradation; completion-only masking is
  the most important follow-up.
  \item \textbf{Single seed for some conditions.} Conditions i, ii, v are single
  runs (i and v are deterministic; gold is single-seed). The gold-vs-zero-shot
  effect is large and significant, but multi-seed gold would tighten it.
  \item \textbf{Approximate MATH grading.} Our matcher is string/numeric, not a
  CAS; reported MATH numbers are conservative lower bounds ($\approx$7\%
  false-negative rate in our audit).
  \item \textbf{Latency hardware.} Latency was measured on A100, not the T4 in our
  proposal; absolute numbers are hardware-specific.
  \item \textbf{Possible GSM8K contamination.} The 1B model may have seen GSM8K
  during pretraining; we mitigate by treating MATH as the genuine-generalization
  signal, where we see no gain.
\end{itemize}

\section{Conclusion and Future Work}
Sequence-level CoT distillation from a small open 8B teacher meaningfully improves a
1B student on GSM8K ($38.1\%\rightarrow46.9\%$), and teacher-generated reasoning
clearly beats the dataset's own gold solutions---both overwhelmingly significant.
Equally informative are our negative findings: gold-CoT fine-tuning hurts, more
teacher traces hurt, and gains do not transfer to MATH. The practical payoff of the
small student is memory ($5.2\times$ lower VRAM), not latency, because the student
inherits the teacher's verbosity. Future work: (1) completion-only loss masking
(highest priority); (2) length-controlled or quality-filtered trace selection to
address the verbosity tax and the counter-productive third trace; (3) a stronger
math teacher (e.g., Qwen2.5-Math-7B) and a CAS-based MATH grader; (4) mixing in a
small amount of OOD data to test whether transfer can be induced.

\begin{thebibliography}{9}
\bibitem{hinton2015} G.~Hinton, O.~Vinyals, J.~Dean. Distilling the Knowledge in a Neural Network. \emph{NeurIPS Workshop}, 2015.
\bibitem{minillm} Y.~Gu, L.~Dong, F.~Wei, M.~Huang. MiniLLM: Knowledge Distillation of Large Language Models. \emph{ICLR}, 2024.
\bibitem{distillm} J.~Ko, S.~Kim, et al. DistiLLM: Towards Streamlined Distillation for Large Language Models. \emph{ICML}, 2024.
\bibitem{deepseekr1} DeepSeek-AI. DeepSeek-R1: Incentivizing Reasoning Capability in LLMs via Reinforcement Learning. 2025.
\bibitem{magister2023} L.~C.~Magister, J.~Mallinson, J.~Adamek, et al. Teaching Small Language Models to Reason. \emph{ACL}, 2023.
\bibitem{lora} E.~Hu, Y.~Shen, et al. LoRA: Low-Rank Adaptation of Large Language Models. \emph{ICLR}, 2022.
\bibitem{qlora} T.~Dettmers, A.~Pagnoni, A.~Holtzman, L.~Zettlemoyer. QLoRA: Efficient Finetuning of Quantized LLMs. \emph{NeurIPS}, 2023.
\bibitem{gsm8k} K.~Cobbe, V.~Kosaraju, et al. Training Verifiers to Solve Math Word Problems. \emph{arXiv:2110.14168}, 2021.
\bibitem{math} D.~Hendrycks, C.~Burns, et al. Measuring Mathematical Problem Solving With the MATH Dataset. \emph{NeurIPS D\&B}, 2021.
\end{thebibliography}

\appendix
\section{Reproducibility}
Code: GitHub repository (\texttt{mathdistill} package + Colab notebook). Pipeline
stages are config-driven (YAML) and resumable. Released artifacts: cleaned
teacher-CoT dataset and best LoRA adapter (Hugging Face). Key settings: GSM8K
\texttt{main}; teacher 4-bit NF4; QLoRA rank 16 / $\alpha$ 32 / 3 epochs / lr
$2\times10^{-4}$; greedy evaluation; 3 seeds for trained conditions.

\section{Qualitative transcripts}
\label{app:qual}
\emph{Paste the four cases from the notebook's qualitative analysis here:
(i) ice-cream win, (ii) vacuum-cleaner win, (iii) overtime-pay loss, (iv) Janet's
ducks style comparison.}

\section{Use of AI assistants}
Per course policy: we used an AI coding assistant to help implement the pipeline
code, debug Colab/TRL environment issues, and draft this report; all experiments,
results, and analysis decisions are our own, and all reported numbers come from our
runs. \emph{(Adjust to reflect your exact usage before submission.)}

\end{document}
